\section{What the models know about masked firms}\label{sec:memorization}

Masking removes names, but news is distinctive: a drug approval, a store count or a guidance range can point to
one company. We measure how often a model recognizes the company behind a masked state. Each of \res{mem-n}
masked val-A states is shown with four candidate companies from the same sector and size tier, the true one
and three distractors chosen to match its news activity, and the model picks one. For model $m$,
\begin{equation}\label{eq:memorization}
\mathrm{Acc}_m = \frac{1}{n}\sum_{i=1}^{n} \mathbf{1}\bigl[\hat c_{i} = c_{i}\bigr],
\qquad
\text{excess}_m = \mathrm{Acc}_m - p_0,
\end{equation}
where $c_i$ is the true company, $\hat c_i$ the model's choice and $p_0 = \res{mem-baseline-f2}$ the matched
baseline: the accuracy of a guess that uses no reading, only the candidates' news activity. Intervals are
Wilson score intervals \citep{wilson1927}. The matched baseline matters: with distractors drawn at random, a
guess based on news activity alone already scores \res{mem-baseline-na06}, because the true company always has
news.

The commercial decision models recognize the company far more often than the baseline
(Figure~\ref{fig:memorization}): from \res{mem-min} of states (mercury-decide) to \res{mem-max} (d1-ultra), and
\res{mem-span} and \res{mem-drex} for the other two comparators. The open-weights Laya model, the base of
Laya-Fin, does so in \res{mem-laya}, close to the baseline. All of the commercial models therefore fail the
pre-registered screen of 40\%, and every model in the figure was released after the sample ended: Laya on
\res{rel-laya} and the commercial models between \res{rel-drex} and \res{rel-span}. Their training data are
undisclosed, so we cannot rule out that they saw text from the test period.

\begin{figure}[tbp]
\centering
\includegraphics{\figdir/fig7_memorisation.pdf}
\caption{\textbf{Commercial decision models identify the company behind masked news in \res{mem-min} to
\res{mem-max} of states, against a matched baseline of \res{mem-baseline}; the open-weights Laya model does so in
\res{mem-laya}.} Company identification on masked val-A states with activity-matched distractors: the share of
states in which each model names the company among four candidates, sorted, with its Wilson 95\% interval,
against the matched baseline (a guess that uses no reading; solid line) and the pre-registered screening
threshold (40\%; dashed line). The three commercial models carried to the test are in the comparator colour,
the other commercial models and Laya zero-shot in gray. Every model was released after the sample ended.
Identifying a company is not the same as remembering its returns. $^{\ddagger}$span-02: see
Figure~\ref{fig:validation-test}.}\label{fig:memorization}
\end{figure}

Recognition is not the same as remembering returns, and two results suggest that it did not drive the
commercial models' signal: unmasking the states does not raise d1-ultra's IC (Section~\ref{sec:priced}), and
their ICs fell from validation to test as the other signals' did, which is not what an advantage from having
seen the test period would produce. Neither is a test of memorization. The result is, however, a warning for
backtests of language models on historical news: masking hides names, not companies. The pre-registered forward
test (Section~\ref{sec:limits}) asks the question directly, on news published after the models were released.
